\documentclass[11pt]{article}
\usepackage[a4paper,margin=1in]{geometry}
\usepackage{amsmath,amssymb,amsthm,mathtools}
\usepackage{graphicx}
\usepackage[hidelinks]{hyperref}
\usepackage{xcolor}

\newtheorem{theorem}{Theorem}
\newtheorem{lemma}[theorem]{Lemma}
\newcommand{\T}{\mathbb T}
\newcommand{\Irr}{\operatorname{Irr}}
\newcommand{\1}{\mathbf 1}
\newcommand{\norm}[1]{\lVert #1\rVert}

\title{The sharp Rider gap for split compact metabelian groups}
\author{Candidate partial solution for arXiv:1110.6683\\
\texttt{agent\_lane\_02} \; (GPT5.6)}
\date{11 August 2026}

\begin{document}
\maketitle

\begin{center}
\fbox{\parbox{0.91\textwidth}{\textbf{Status: candidate partial result, likely valid.}
The theorem below proves the conjecturally sharp universal constant for every compact
semidirect product of two abelian groups.  It does not settle the source question for
arbitrary compact groups.}}
\end{center}

\section{Source question}

The source is M.~Alaghmandan, Y.~Choi and E.~Samei,
\emph{$ZL$-amenability and characters for the restricted direct products of finite groups},
arXiv:1110.6683v4 (2013)~\cite{ACS}.  Theorem~2.3 recalls Rider's estimate for a finite
linear combination $\psi$ of irreducible characters on a compact group: if
$\psi*\psi=\psi$ and $\norm{\psi}_1>1$, then $\norm{\psi}_1\geq301/300$.
Immediately afterwards, Remark~2.4 asks:

\begin{quote}
``We also note that $301/300$ is not sharp; it appears to be an open problem to
determine the best constants here.''
\end{quote}

The passage occurs on physical PDF page~4 (source Section~2.2):

\begin{center}
\includegraphics[width=0.96\textwidth]{figures/open_problem_crop.png}
\end{center}

We use normalized Haar measures.  This is equivalent to the source formulation by
the rescaling noted in its Remark~2.4.  For an abelian compact group, Saeki proved the
sharp alternative
\[
  \norm f_1=1\qquad\hbox{or}\qquad
  \norm f_1\geq c_{\rm S}:=\frac{1+\sqrt2}{2}
\]
for nonzero idempotent measures~\cite{Saeki}.  The question is whether the same
constant works in the noncommutative setting.

\section{Proof intuition}

Let $G=A\rtimes H$, where both $A$ and $H$ are compact abelian.  The irreducible
representations of $G$ are the familiar little-group representations.  An $H$-orbit
$\mathcal O\subset\widehat A$, together with a character of its stabilizer, labels one
irreducible representation.  The key observation is that its minimal central
idempotent $d_\pi\chi_\pi$ is, point by point on the underlying space $A\times H$,
an ordinary sum of distinct characters of the \emph{abelian} direct product
$A\times H$.  The associated character rectangles are disjoint and partition
$\widehat A\times\widehat H$.

Consequently, every finite-spectrum central idempotent on the nonabelian semidirect
product has exactly the same values, and hence exactly the same $L^1$ norm, as an
idempotent trigonometric polynomial on the abelian direct product.  Saeki's sharp
abelian theorem applies without loss in the constant.  This transference is special
to the split abelian-by-abelian structure; the irreducible dual of a general compact
group does not admit the required disjoint character rectangles.

\section{The transference identity}

Let $A,H$ be compact abelian groups, let $\alpha:H\to\operatorname{Aut}(A)$ be a
continuous action, and put $G=A\rtimes_\alpha H$.  We identify the underlying
probability spaces of $G$ and the abelian direct product $B=A\times H$; in both cases
normalized Haar measure is $m_A\times m_H$.

For $\gamma\in\widehat A$, write
\[
 \mathcal O_\gamma=H\cdot\gamma,
 \qquad H_\gamma=\{h\in H:h\cdot\gamma=\gamma\},
 \qquad d_\gamma=[H:H_\gamma].
\]
The orbit is finite: it is a compact subset of the discrete group $\widehat A$.
For $\eta\in\widehat{H_\gamma}$, let
\[
 \pi_{\gamma,\eta}
 =\operatorname{Ind}_{A\rtimes H_\gamma}^{A\rtimes H}(\widetilde\gamma\otimes\eta),
 \qquad \widetilde\gamma(a,h)=\gamma(a).
\]
This is well-defined because $h$ fixes $\gamma$ on $H_\gamma$.

\begin{lemma}[Rectangle lemma]\label{lem:rectangle}
Choose one $\gamma$ from each $H$-orbit in $\widehat A$.  The representations
$\pi_{\gamma,\eta}$, with $\eta\in\widehat{H_\gamma}$, form $\Irr(G)$.
For
\[
 X_{\gamma,\eta}:=\mathcal O_\gamma\times
 \{\xi\in\widehat H:\xi|_{H_\gamma}=\eta\}
 \subseteq\widehat A\times\widehat H,
\]
the sets $X_{\gamma,\eta}$ are finite, pairwise disjoint, and partition
$\widehat A\times\widehat H$.  Moreover, for every $(a,h)\in A\times H$,
\begin{equation}\label{eq:rectangle}
 d_{\pi_{\gamma,\eta}}\chi_{\pi_{\gamma,\eta}}(a,h)
 =\sum_{(\delta,\xi)\in X_{\gamma,\eta}}\delta(a)\xi(h).
\end{equation}
\end{lemma}

\begin{proof}
We first recall the short little-group argument.  If $\pi$ is irreducible, its
restriction to $A$ splits into finitely many weight spaces.  Normality of $A$ makes
$H$ permute the weights, and irreducibility makes this action transitive.  For a
weight $\gamma$, the abelian stabilizer $H_\gamma$ acts diagonally on its weight
space.  The $G$-span of any common eigenline is invariant; irreducibility therefore
forces the weight space to be one-dimensional and shows that $\pi$ is precisely one
of the displayed induced representations.  This also proves completeness and the
usual uniqueness up to the $H$-orbit of $\gamma$.

The induced character formula gives, with $d=d_\gamma$,
\begin{equation}\label{eq:induced}
 \chi_{\pi_{\gamma,\eta}}(a,h)
 =\begin{cases}
   \eta(h)\displaystyle\sum_{\delta\in\mathcal O_\gamma}\delta(a),&h\in H_\gamma,\\[5pt]
   0,&h\notin H_\gamma.
  \end{cases}
\end{equation}
Indeed, $A\rtimes H_\gamma$ has index $d$, and conjugation by the $H$-coset
representatives produces exactly the weights in $\mathcal O_\gamma$.

Restriction $\widehat H\to\widehat{H_\gamma}$ is onto by Pontryagin duality.
The fiber over $\eta$ is a coset of the annihilator $H_\gamma^\perp$, which has
$d$ elements.  Finite-character orthogonality therefore gives
\[
 \sum_{\substack{\xi\in\widehat H\\ \xi|_{H_\gamma}=\eta}}\xi(h)
 =\begin{cases}d\eta(h),&h\in H_\gamma,\\0,&h\notin H_\gamma.\end{cases}
\]
Multiplying this identity by the orbit sum in \eqref{eq:induced} proves
\eqref{eq:rectangle}, since $d_{\pi_{\gamma,\eta}}=d$.

For a fixed orbit, the restriction fibers partition $\widehat H$; distinct
$H$-orbits in $\widehat A$ are disjoint.  Hence the rectangles are pairwise
disjoint and partition $\widehat A\times\widehat H$.
\end{proof}

\begin{theorem}[Sharp Rider gap for split compact metabelian groups]\label{thm:main}
Let $G=A\rtimes H$, where $A$ and $H$ are compact abelian groups.  Let $\psi$ be a
finite linear combination of irreducible characters of $G$ satisfying
$\psi*\psi=\psi$ with respect to normalized Haar measure.  Then
\[
 \norm{\psi}_1\in\{0,1\}\quad\text{or}\quad
 \norm{\psi}_1\geq\frac{1+\sqrt2}{2}.
\]
The positive gap constant is sharp even among nonabelian finite groups of this form.
\end{theorem}

\begin{proof}
Peter--Weyl orthogonality says that a finite-spectrum central convolution
idempotent has the unique form
\[
 \psi=\sum_{\pi\in S}d_\pi\chi_\pi
\]
for a finite set $S\subseteq\Irr(G)$.  Let
\[
 E=\bigsqcup_{\pi_{\gamma,\eta}\in S}X_{\gamma,\eta}
 \subseteq\widehat B,
 \qquad
 F(a,h)=\sum_{(\delta,\xi)\in E}\delta(a)\xi(h).
\]
The union is disjoint by Lemma~\ref{lem:rectangle}.  Thus $F$ is an idempotent
trigonometric polynomial on the compact abelian group $B=A\times H$.
Equation~\eqref{eq:rectangle}, summed over $S$, gives the \emph{pointwise} identity
\[
 \psi(a,h)=F(a,h)\qquad((a,h)\in A\times H),
\]
where the left side uses the semidirect multiplication and the right side the direct
product multiplication.  The normalized Haar measures coincide on the underlying
space, so $\norm\psi_1=\norm F_1$.  Saeki's sharp abelian norm theorem~\cite{Saeki}
now yields the stated alternative.

For sharpness take $G=D_{16}=C_8\rtimes C_2$, with inversion action.  Let $\pi$ be
the two-dimensional irreducible representation associated to the orbit
$\{1,-1\}\subset\widehat{C_8}$.  For $\psi=2\chi_\pi$, one has
\[
 \psi(r^k)=4\cos(\pi k/4),\qquad \psi(sr^k)=0,
\]
and hence
\[
 \norm\psi_1
 =\frac1{16}\sum_{k=0}^{7}\left|4\cos(\pi k/4)\right|
 =\frac{1+\sqrt2}{2}.
\]
\end{proof}

\section{Verification and limitations}

The reusable script \texttt{code/verifier.py} checks the rectangle identity and
partition property for 3,160 little-group blocks in cyclic semidirect products
$C_n\rtimes C_m$ with $2\le n\le14$, $2\le m\le6$, over every admissible cyclic
action in that range.  It also verifies the exact $D_{16}$ equality and exhaustively
checks all central idempotents of $D_{2n}$ for $3\le n\le30$.  These computations are
sanity checks, not a substitute for Lemma~\ref{lem:rectangle}.  Details and the exact
command appear in \texttt{VERIFIER\_REPORT.md}.

The original question ranges over all compact groups.  The proof here requires a
\emph{split} extension with both normal factor and complement abelian.  For a general
compact group, irreducible representations do not yield a disjoint partition of an
abelian dual into character rectangles, so the transference step is unavailable.

\paragraph{Novelty check.}
A bounded search on 11 August 2026 used the exact source phrase, ``Rider'' with
``best constant'', ``central idempotent'' with ``semidirect product'', ``metabelian'',
and ``abelian normal subgroup'', plus forward citations visible from the source.
The closest structural paper found was Hauenschild's 1979 characterization of central
idempotent measures for groups with an abelian normal subgroup of finite index
\cite{Hauenschild}; its theorem is phrased through hypercoset rings and the inspected
six-page text does not state this norm-preserving direct-product transference or a sharp
gap.  A 2023 thesis still records the sharp finite-group bound as open and reports
computations supporting $c_{\rm S}$~\cite{Sawatzky}.  Thus this packet claims a new
partial theorem with \emph{moderate}, not definitive, novelty confidence; it may also
be a previously unrecorded corollary of older structure theory.

\paragraph{Human-review recommendation.}
Check first the compact little-group classification and the restriction-fiber
orthogonality in Lemma~\ref{lem:rectangle}.  Then check that the rectangles remain
disjoint when several irreducibles are summed.  Once those points are accepted, the
reduction to Saeki is exact and the sharpness computation is elementary.

\begin{thebibliography}{9}
\bibitem{ACS}
M.~Alaghmandan, Y.~Choi and E.~Samei,
\emph{$ZL$-amenability and characters for the restricted direct products of finite groups},
arXiv:1110.6683v4, 2013.

\bibitem{Saeki}
S.~Saeki, \emph{On norms of idempotent measures},
Proc. Amer. Math. Soc. \textbf{19} (1968), 600--602,
doi:10.1090/S0002-9939-1968-0225102-6.

\bibitem{Hauenschild}
W.~Hauenschild, \emph{Zentrale idempotente Ma\ss e auf lokalkompakten Gruppen mit
Abelschem Normalteiler von endlichem Index}, Manuscripta Math. \textbf{28} (1979),
405--410.

\bibitem{Sawatzky}
J.~Sawatzky, \emph{Mind the GAP: Amenability}, Master's thesis,
University of Waterloo, 2023, Chapter~6, Conjecture~6.2.3.
\end{thebibliography}

\end{document}
